% GENERATED FILE. Edit canonical lessons, then run npm run generate:books.
% canonical-source-hash: fnv1a64:a3b8190660a8bb69
% canonical-lessons: RU-C495-beremennost, RU-C495-roddom, RU-C495-dekret, RU-C495-kolyaska, RU-C495-polzat

\chapter{Expecting a Baby}
\label{ch:ru-expecting-a-baby}
\hlchaptermodality{495}

\begin{chapteropening}
\textbf{By the end of this chapter:} \emph{I can talk about pregnancy and a new baby.}

Complete the last lesson of chapter 495: I can talk about pregnancy and a new baby.
\end{chapteropening}

\section[\texorpdfstring{\ru{беременность} (berémennost)}{беременность (berémennost)}]{\ru{беременность} (berémennost) — pregnancy}

\label{lesson:RU-C495-beremennost}



\begin{quote}
Before the new one: say the Russian for to be unfaithful, then the Russian for a parting.
\end{quote}

\subsection*{You'll want to know: \ru{беременность}}
\textbf{\ru{беременность}} (\emph{berémennost}) — \textquotedblleft{}pregnancy\textquotedblright{}.

\begin{quote}
\ru{Во} \ru{время} \ru{беременности} \ru{она} \ru{много} \ru{гуляла}.
\end{quote}

\begin{quote}
\emph{Vo vrémya berémennosti oná mnógo gulyála.}
\end{quote}

\begin{quote}
During her pregnancy she walked a lot.
\end{quote}

Feminine.

\begin{cousinweb}
\emph{\ru{беременность}} is built on \emph{\ru{беременная}} 'pregnant', from the older \emph{\ru{беремя}} 'a load, an armful', usually linked to \emph{\ru{брать}} 'to take'; EtymDB has no entry.
\end{cousinweb}

\begin{grammarlens}[title={What you've built}]
One more word for telling, explaining and describing at B1.
\end{grammarlens}

\subsection*{Your turn}
\begin{itemize}
\raggedright
  \item \emph{Say it:} \emph{berémennost}
  \item \emph{Say it:} \emph{berémennost}, once more
  \item \emph{Say it:} a sentence of your own with \emph{berémennost}
  \item \emph{Recall:} say the Russian for to be unfaithful, then the Russian for a parting, then say \emph{berémennost} again
\end{itemize}

\begin{tcolorbox}[breakable,colback=teal!4,colframe=teal!35!black,title={Before you move on}]
What does \ru{беременность} mean? (\textquotedblleft{}Pregnancy\textquotedblright{}.) Say it once more.
\end{tcolorbox}
\section[\texorpdfstring{\ru{роддом} (roddóm)}{роддом (roddóm)}]{\ru{роддом} (roddóm) — a maternity hospital}

\label{lesson:RU-C495-roddom}



\begin{quote}
Before the new one: say the Russian for a parting, then the Russian for pregnancy.
\end{quote}

\subsection*{You'll want to know: \ru{роддом}}
\textbf{\ru{роддом}} (\emph{roddóm}) — \textquotedblleft{}a maternity hospital\textquotedblright{}.

\begin{quote}
\ru{Папа} \ru{встретил} \ru{маму} \ru{и} \ru{сестрёнку} \ru{у} \ru{роддома}.
\end{quote}

\begin{quote}
\emph{Pápa vstrétil mámu i sestryónku u roddóma.}
\end{quote}

\begin{quote}
Dad met Mum and my baby sister at the maternity hospital.
\end{quote}

Masculine. Short for \emph{\ru{родильный} \ru{дом}}.

\begin{cousinweb}
\emph{\ru{роддом}} is a clipped compound of \emph{\ru{родильный}} 'birth-', from \emph{\ru{родить}} 'to give birth', and \emph{\ru{дом}} 'house'; EtymDB has no entry.
\end{cousinweb}

\begin{grammarlens}[title={What you've built}]
One more word for telling, explaining and describing at B1.
\end{grammarlens}

\subsection*{Your turn}
\begin{itemize}
\raggedright
  \item \emph{Say it:} \emph{roddóm}
  \item \emph{Say it:} \emph{roddóm}, once more
  \item \emph{Say it:} a sentence of your own with \emph{roddóm}
  \item \emph{Recall:} say the Russian for a parting, then the Russian for pregnancy, then say \emph{roddóm} again
\end{itemize}

\begin{tcolorbox}[breakable,colback=teal!4,colframe=teal!35!black,title={Before you move on}]
What does \ru{роддом} mean? (\textquotedblleft{}A maternity hospital\textquotedblright{}.) Say it once more.
\end{tcolorbox}
\section[\texorpdfstring{\ru{декрет} (dekrét)}{декрет (dekrét)}]{\ru{декрет} (dekrét) — maternity leave; a decree}

\label{lesson:RU-C495-dekret}



\begin{quote}
Before the new one: say the Russian for pregnancy, then the Russian for a maternity hospital.
\end{quote}

\subsection*{You'll want to know: \ru{декрет}}
\textbf{\ru{декрет}} (\emph{dekrét}) — \textquotedblleft{}maternity leave; a decree\textquotedblright{}.

\begin{quote}
\ru{После} \ru{рождения} \ru{сына} \ru{она} \ru{год} \ru{была} \ru{в} \ru{декрете}.
\end{quote}

\begin{quote}
\emph{Pósle rozhdéniya sýna oná god bylá v dekréte.}
\end{quote}

\begin{quote}
After her son was born she was on maternity leave for a year.
\end{quote}

Masculine. Colloquial for \emph{\ru{декретный} \ru{отпуск}}. Its older sense is a decree.

\begin{cousinweb}
\emph{\ru{декрет}} was borrowed from Latin \emph{decretum} 'a decree' (EtymDB), from \emph{dēcernere} 'to decide' (Lewis \& Short).
\end{cousinweb}

\begin{grammarlens}[title={What you've built}]
One more word for telling, explaining and describing at B1.
\end{grammarlens}

\subsection*{Your turn}
\begin{itemize}
\raggedright
  \item \emph{Say it:} \emph{dekrét}
  \item \emph{Say it:} \emph{dekrét}, once more
  \item \emph{Say it:} a sentence of your own with \emph{dekrét}
  \item \emph{Recall:} say the Russian for pregnancy, then the Russian for a maternity hospital, then say \emph{dekrét} again
\end{itemize}

\begin{tcolorbox}[breakable,colback=teal!4,colframe=teal!35!black,title={Before you move on}]
What does \ru{декрет} mean? (\textquotedblleft{}Maternity leave; a decree\textquotedblright{}.) Say it once more.
\end{tcolorbox}
\section[\texorpdfstring{\ru{коляска} (kolyáska)}{коляска (kolyáska)}]{\ru{коляска} (kolyáska) — a pram; a wheelchair}

\label{lesson:RU-C495-kolyaska}



\begin{quote}
Before the new one: say the Russian for a maternity hospital, then the Russian for maternity leave.
\end{quote}

\subsection*{You'll want to know: \ru{коляска}}
\textbf{\ru{коляска}} (\emph{kolyáska}) — \textquotedblleft{}a pram; a wheelchair\textquotedblright{}.

\begin{quote}
\ru{Каждый} \ru{день} \ru{мы} \ru{гуляем} \ru{с} \ru{коляской} \ru{в} \ru{парке}.
\end{quote}

\begin{quote}
\emph{Kázhdyy den my gulyáyem s kolyáskoy v párke.}
\end{quote}

\begin{quote}
Every day we go for a walk in the park with the pram.
\end{quote}

Feminine. \emph{\ru{Детская} \ru{коляска}} — a pram; \emph{\ru{инвалидная} \ru{коляска}} — a wheelchair.

\begin{cousinweb}
EtymDB lists \emph{\ru{коляска}} without a source; it is usually traced through Polish \emph{kolaska} to the root of \emph{\ru{колесо}} 'wheel'.
\end{cousinweb}

\begin{grammarlens}[title={What you've built}]
One more word for telling, explaining and describing at B1.
\end{grammarlens}

\subsection*{Your turn}
\begin{itemize}
\raggedright
  \item \emph{Say it:} \emph{kolyáska}
  \item \emph{Say it:} \emph{kolyáska}, once more
  \item \emph{Say it:} a sentence of your own with \emph{kolyáska}
  \item \emph{Recall:} say the Russian for a maternity hospital, then the Russian for maternity leave, then say \emph{kolyáska} again
\end{itemize}

\begin{tcolorbox}[breakable,colback=teal!4,colframe=teal!35!black,title={Before you move on}]
What does \ru{коляска} mean? (\textquotedblleft{}A pram; a wheelchair\textquotedblright{}.) Say it once more.
\end{tcolorbox}
\section[\texorpdfstring{\ru{ползать} (pólzat)}{ползать (pólzat)}]{\ru{ползать} (pólzat) — to crawl}

\label{lesson:RU-C495-polzat}



\begin{quote}
Before the new one: say the Russian for maternity leave, then the Russian for a pram.
\end{quote}

\subsection*{You'll want to know: \ru{ползать}}
\textbf{\ru{ползать}} (\emph{pólzat}) — \textquotedblleft{}to crawl\textquotedblright{}.

\begin{quote}
\ru{Наш} \ru{сын} \ru{уже} \ru{ползает} \ru{по} \ru{всей} \ru{квартире}.
\end{quote}

\begin{quote}
\emph{Nash syn uzhé pólzayet po vsey kvartíre.}
\end{quote}

\begin{quote}
Our son already crawls all over the flat.
\end{quote}

Imperfective, in many directions; the one-way verb is \emph{\ru{ползти}}.

\begin{cousinweb}
\emph{\ru{ползать}} is inherited from Proto-Slavic \emph{p\ru{ь}lzati} (EtymDB).
\end{cousinweb}

\begin{grammarlens}[title={What you've built}]
One more word for telling, explaining and describing at B1.
\end{grammarlens}

\subsection*{Your turn}
\begin{itemize}
\raggedright
  \item \emph{Say it:} \emph{pólzat}
  \item \emph{Say it:} \emph{pólzat}, once more
  \item \emph{Say it:} a sentence of your own with \emph{pólzat}
  \item \emph{Recall:} say the Russian for maternity leave, then the Russian for a pram, then say \emph{pólzat} again
\end{itemize}

\begin{tcolorbox}[breakable,colback=teal!4,colframe=teal!35!black,title={Before you move on}]
What does \ru{ползать} mean? (\textquotedblleft{}To crawl\textquotedblright{}.) Say it once more.
\end{tcolorbox}
